\documentclass[]{IEEEtran}
% some very useful LaTeX packages include:
%\usepackage{cite}      
\usepackage{graphicx}   
\usepackage{subfigure} 
\usepackage{url}       
\usepackage{amsmath}    
\usepackage{caption2}
% Your document starts here!
\begin{document}

% Define document title and author
	\title{Weekly Report}
	\author{Adviser: Prof. Yang Wen \\Student: Cheng Wensheng\\ Period: 2018.3.19- 3.24
	}
	\markboth{Visual Information Processing Group}{}
	\maketitle

% Write abstract here
\begin{abstract}
	This week I mainly put my effort on training PSPNet and RefineNet model on our dataset.
\end{abstract}

% Each section begins with a \section{title} command
\section{Model Training}
	% \PARstart{}{} creates a tall first letter for this first paragraph
	\PARstart{W}{e} tried most of the state-of-the-art models, including DeepLab V3, DenseNet, PSPNet, RefineNet and GCN model.
	\begin{itemize}
		\item We tried DeepLab V3 and DenseNet last week. However, neither performed well on our dataset. We have no choice but to try more.
		\item After a long period of training, the result showed these models didn't perform well as I expected. Although they get competitive results on PASCAL VOC dataset, such as \textbf{80$\%$ mIoU}, they only achieve \textbf{50-60$\%$ mIoU} on our dataset in turn.
	\end{itemize}

% Main Part
\section{Result Analysis}
	% LaTeX takes complete care of your document layout ...
	We discussed about our result, then raised following reasons for poor accuracy.
	\begin{itemize}
		\item Obviously, the mere difference is dataset. We checked up them and compared carefully. One major drawback is the fuzzy label of our dataset. For optical remote sensing images, some places are shaded by trees or high buildings. Yet for label images, these regions are labeled as primary class. Therefore, they would pollute the whole dataset and make the model confused.   
		\item As for our result, the lowest accuracy is of road class. Since the color of road is very close to some roofs, and they are both flat, it's hard for the model to distinguish road from these buildings.  
		\item The test result of these models in sequences are Fig.~\ref{fig:fcd_gt},  Fig.~\ref{fig:fcd_pr}.
	\end{itemize}
\newpage
\begin{figure}[!hbt]
%		 Center the figure.
		\vspace{2.5cm}
%		\hspace{50cm}
		\begin{center}
			\begin{minipage}[t]{0.21\textwidth}
				\includegraphics[width=\columnwidth]{dl3_SAT01-00008_gt}
				%		 Create a subtitle for the figure.
				\caption{Ground truth image DeepLab V3}
				\label{fig:dlv3_gt}
			\end{minipage}
		    \hspace{0.5cm}
			\begin{minipage}[t]{0.21\textwidth}
				\includegraphics[width=\columnwidth]{dl3_SAT01-00008_pred}
				%Create a subtitle for the figure.
				\caption{Prediction image DeepLab V3}
				\label{fig:dlv3_pr}
			\end{minipage}
		\end{center}
	\end{figure}

\begin{figure}[!hbt]
	%		 Center the figure.
%	\vspace{1.8cm}
	\begin{center}
		\vspace{-0.3cm}
		\begin{minipage}[t]{0.21\textwidth}
			\includegraphics[width=\columnwidth]{fc_SAT01-00008_gt}
			%		 Create a subtitle for the figure.
			\caption{Ground truth image FCDenseNet}
			\label{fig:fcd_gt}
		\end{minipage}
		\hspace{0.5cm}
		\begin{minipage}[t]{0.21\textwidth}
			\includegraphics[width=\columnwidth]{fc_SAT01-00008_pred}
			%Create a subtitle for the figure.
			\caption{Prediction image FCDenseNet}
			\label{fig:fcd_pr}
		\end{minipage}
	\end{center}
\end{figure}

% Now we need a bibliography:
%\begin{thebibliography}{5}
%
%	%Each item starts with a \bibitem{reference} command and the details thereafter.
%	\bibitem{HOP96} % Transaction paper
%	J.~Hagenauer, E.~Offer, and L.~Papke. Iterative decoding of binary block
%	and convolutional codes. {\em IEEE Trans. Inform. Theory},
%	vol.~42, no.~2, pp.~429–-445, Mar. 1996.
%
%	\bibitem{MJH06} % Conference paper
%	T.~Mayer, H.~Jenkac, and J.~Hagenauer. Turbo base-station cooperation for intercell interference cancellation. {\em IEEE Int. Conf. Commun. (ICC)}, Istanbul, Turkey, pp.~356--361, June 2006.
%
%	\bibitem{Proakis} % Book
%	J.~G.~Proakis. {\em Digital Communications}. McGraw-Hill Book Co.,
%	New York, USA, 3rd edition, 1995.
%
%	\bibitem{talk} % Web document
%	F.~R.~Kschischang. Giving a talk: Guidelines for the Preparation and Presentation of Technical Seminars.
%	\url{http://www.comm.toronto.edu/frank/guide/guide.pdf}.
%
%	\bibitem{5}
%	IEEE Transactions \LaTeX and Microsoft Word Style Files.
%	\url{http://www.ieee.org/web/publications/authors/transjnl/index.html}
%
%\end{thebibliography}

% Your document ends here!
\end{document}